% Copyright 2026 Torbjörn E. M. Nordling
% SPDX-License-Identifier: Apache-2.0
% frame-git-setup.tex — Before the hand-in: accounts and tools (22:30–23:15)
% Students have no GitHub account, no git, and do not know what a fork is
% (user, 2026-11-12). Steps 1 to 4 of the student README "Step by step";
% screenshots are the README's own, in the repository's assets/image/.
\begin{frame}[fragile]{Before the hand-in: accounts and tools}{Step by step, with screenshots, in the README of the repository}
  \scriptsize
  \begin{columns}[T]
    \column{0.6\textwidth}
    \begin{enumerate}\setlength{\itemsep}{1pt}
      \item \textbf{Everyone:} create a GitHub account at \texttt{github.com/signup}.
      \item \textbf{Windows only:} install git from \texttt{git-scm.com}, all defaults; use \alert{Git Bash}. Everyone, once:
\begin{verbatim}
git config --global user.name "<name>"
git config --global user.email "<email>"
\end{verbatim}
      \item \textbf{One member:} \alert{Fork} the course repository, your group's copy; then \textbf{Settings}, \textbf{Collaborators}, \textbf{Add people}.
      \item \textbf{Everyone:} copy the fork to your computer:
\begin{verbatim}
git clone <fork address>
cd <repository folder>
git remote add upstream <course address>
\end{verbatim}
      First push: macOS and Linux use a \alert{token} as password.
      \item \textbf{iPad:} Safari on github.com; analyse on a laptop.
    \end{enumerate}
    \column{0.37\textwidth}
    \centering
    \includegraphics[width=\linewidth]{github-fork-create.jpg}\\[2pt]
    \includegraphics[width=\linewidth]{github-fork-collaborators.jpg}
  \end{columns}
\end{frame}
